\chapter{Getting started}
\label{ch:quickstart}

This chapter walks through installing the system, the interactive workspace, and two complete worked examples.  It uses the vocabulary of Part~\ref{part:system}; a reader who wants to run something first may start here.

%% (moved from ch-intro.tex)
\label{sec:quickstart-chap}

\section{Installation and startup}

\subsection{Requirements}

MIT/GNU~Scheme~12.1, and GNU~Emacs~\cite{Emacs} (for the interactive workspace).
A \texttt{Makefile} at the root of the tree wraps the steps below
(\texttt{make help} lists them).
No other dependencies.

\subsection{Directory layout}

The prover lives in \texttt{\textasciitilde/prover/}: the executable
script \texttt{prover} and the loader \texttt{load.scm} are the only
two files you need to know about to get started.  See
Appendix~\ref{app:source} for a per-file overview.

\subsection{Compile the tree first}
\label{sec:compile-first}

Compiled binaries (\texttt{.com}, \texttt{.bin}) are excluded from both the
repository and the distributed archive, so a freshly cloned or unpacked tree is
\emph{interpreted}.  It will run correctly, and nothing in its output will say
that it is doing so; only the clock will.  On a modest machine the difference is
between a library load of seconds and one of minutes, and likewise for the
full test suite.

Compile once, after obtaining the tree.  This includes one load from source:

\begin{VNBsnip}
~/prover/VNB-with-compile    # compile every source file, then launch
\end{VNBsnip}

Thereafter recompile only what has changed.  From a session that has already
loaded \texttt{load.scm}:

\begin{VNBsnip}
(compile-vnb!)               # incremental: skips files whose .com is current
\end{VNBsnip}

After editing a source file, recompile it before running anything.  The
loader compares modification times: a \texttt{.com} older than its
\texttt{.scm} is ignored and the file loads from source.  That is harmless
for a proof script and costly for a core file, and it is the usual
explanation for a session that has become slow.  The question to ask is
whether \emph{that file} is compiled.  A small number of files are compiled
by nobody, deliberately: those that use a top-level macro, which
\texttt{compile-vnb!} lists in \texttt{*vnb-no-compile-files*} and skips.  They
load from source at no measurable cost, their work being done in the compiled
procedures they call.

\section{Starting the checker}

\begin{VNBsnip}
~/prover/prover              # interactive REPL
~/prover/prover proof.scm    # run a proof script and exit
~/prover/prover -i proof.scm # run a script, then stay interactive
\end{VNBsnip}

The interactive session is a standard MIT~Scheme REPL with the proof
checker pre-loaded.
\section{The VNB workspace}
\label{sec:emacs}

The VNB workspace is a styled Emacs session built for one task:
writing proofs.  Start it from the shell:

\begin{VNBsnip}
~/prover/VNB
\end{VNBsnip}

\noindent
A graphical Emacs frame opens with the VNB menu and toolbar already
in place, the prover ready to launch on first use, and a Home
Workspace inviting you to begin.

\begin{remark}
The launcher uses \texttt{emacs -Q}, so it ignores your personal
Emacs init file.  The session is dedicated to the prover; nothing
you do inside it affects your usual Emacs setup.
\end{remark}

\subsection{The Home Workspace}

The first workspace you see.  It is a list of buttons, each with a line of
explanation beside it (\texttt{emacs/vnb-launch.el}, \texttt{vnb-launch--paint-workspace}).
They fall into groups; the list is not exhaustive and the workspace itself is
the authority.

\begin{itemize}
\item \textbf{Starting work.}  \textbf{Start Proof} enters a formula in string
  syntax, switches to the Focus Workspace and begins the proof;
  \textbf{Load Proof Drive} runs a prepared script and stops on the leaf it
  leaves open; \textbf{Calculator} evaluates an arithmetic expression;
  \textbf{Build Structure} defines a new structure and saves it to a file in
  \texttt{structure-library/}.
\item \textbf{Looking things up.}  \textbf{What Is\ldots?}, \textbf{Show Theorems},
  \textbf{Find Theorem}, \textbf{Show PSS}, \textbf{Definitions},
  \textbf{Describe Structure}, \textbf{Structure Manual} and
  \textbf{Notes} show a structure, a number or constant, the installed
  theorems, the Proof Support Set, definitions, a structure card, all
  structures in one document, or a hand-written note.  \textbf{Proof as PDF}
  typesets a stored proof.
\item \textbf{Browsing the library.}  \textbf{Browse Library},
  \textbf{Fingerprint Index}, \textbf{Structure Graph}, \textbf{Library
  (HTML)} and \textbf{Browser Home} open the generated reference pages, some
  of them in a browser.
\item \textbf{Learning and scratch work.}  \textbf{Examples} steps through
  worked example proofs; \textbf{Scratch Workspace} is a Lisp-interaction
  sheet whose \texttt{C-j} sends an expression or region to the prover.
\item \textbf{Session.}  \textbf{Save Session} writes every proof completed
  in the session to one re-loadable script; \textbf{Quit} shuts down the
  prover and exits.
\end{itemize}

\noindent
Most buttons have a single-letter key (for example \texttt{s} starts a proof,
\texttt{t} shows theorems, \texttt{q} quits); the key legend is printed at the
foot of the workspace.  Press \texttt{g} to refresh.

\subsection{The Focus Workspace}

The working surface.  After \textbf{Start Proof} the workspace shows
the goal you are currently focused on, with its assumptions, and, above it, a
short row of buttons (\texttt{vnb-launch--paint-proof-1}):

\begin{itemize}
\item \textbf{What now?} --- asks the copilot what to try on the focused
  goal; each suggested move can be run from there.
\item \textbf{What is\ldots?} --- looks up a structure, number or constant.
\item \textbf{Preamble} --- runs the editable rule file on the focused goal:
  routine steps, then a report of what fired and where it stopped.
\item \textbf{Graph} --- walks the deduction graph one sequent node at a
  time, closed branches included (shown only while a proof is open).
\item \textbf{Undo} --- rolls back the last step (shown only while a proof
  is open).
\end{itemize}

\noindent
The tactics themselves are not buttons.  They run from single-letter keys,
printed at the foot of the workspace in groups: \emph{goal} (\texttt{d} direct
inference, \texttt{u} undo), \emph{hypotheses} (\texttt{a} assume,
\texttt{A} assume all, \texttt{D} decompose a hypothesis, \texttt{e}
separation elimination), \emph{rewriting} (\texttt{m} on the goal,
\texttt{M} on a hypothesis), \emph{citing} (\texttt{t} theorem, \texttt{F}
fact, \texttt{i} instantiate, \texttt{w} witness, \texttt{b} backchain,
\texttt{B} cite a lemma), \emph{closers} (\texttt{p} propositional, \texttt{y}
and \texttt{Y} syllogism), \emph{moving} (\texttt{f} focus, \texttt{o} overview,
\texttt{q} qed) and \emph{workspace} (\texttt{h} home, \texttt{r} scratch pad,
\texttt{S} scratch workspace, \texttt{T} toggle TeX rendering, \texttt{W} save
the proof script, \texttt{g} refresh).  The same moves are in the VNB menu and
can be typed on the scratch sheet.

\begin{remark}
Many more tactics are available in the prover than there are
keys.  The keys cover the most common everyday moves; the
full set is documented in Chapter~\ref{ch:writing-proofs}, and any
tactic can be typed directly in the REPL described below.
\end{remark}

\subsection{The Proof Overview Workspace}

A bird's-eye view of every open goal in the current proof.  Each
goal is rendered as assumptions above a sequent arrow above the
goal formula, with a \textbf{Focus on this} button beside it.
The current focus is marked with a small triangle.  Use this
workspace to navigate large proofs or to remember which goals are
still open.

Press \texttt{f} to return to the Focus Workspace, \texttt{h} for
Home, \texttt{n}/\texttt{p} to jump between goals, and
\texttt{RET} to focus on the goal at point.

\subsection{Error reporting}

When the prover rejects something --- a malformed formula, a tactic
that does not apply --- the error appears in two places at once:

\begin{itemize}
\item the echo area at the bottom of the frame flashes the message
  immediately;
\item a red \texttt{! Error:} panel appears at the top of the
  current workspace and stays there until the next successful
  action clears it.
\end{itemize}

\noindent
No detail is hidden in a side window: whatever the prover said,
you see.

\subsection{Display preferences}

The \textbf{Display} submenu of the VNB menu offers the usual
adjustments: frame size, position, font family and size,
foreground/background/cursor colours.  Any change is saved
automatically to \texttt{\textasciitilde/.vnb-display.el}, so the
next \texttt{VNB} launch reopens the frame exactly the way you
left it.  \textbf{Reset to Defaults} restores the built-in
appearance and removes the saved file.

\subsection{Advanced: the bare REPL}

Sitting behind the styled workspaces is an ordinary MIT~Scheme
REPL.  Open it with \textbf{Show Scratch Pad (advanced)} from the VNB
menu, or skip the workspaces entirely:

\begin{VNBsnip}
~/prover/VNB --no-workspace       # plain comint layout
~/prover/prover                   # no Emacs at all
\end{VNBsnip}

\noindent
The plain Emacs layout is a two-window arrangement: the
\texttt{*VNB*} REPL on the left, the live proof-state display
\texttt{*VNB State*} on the right.  Any Scheme expression you type
at the REPL prompt is evaluated; \texttt{(show)} refreshes the
state display, and short forms like \texttt{(di)} or \texttt{(ass)}
work exactly as they do from the workspace keys.

\paragraph{Scratch sheet of proof commands.}
\texttt{M-x vnb-command-buffer} opens a \texttt{*VNB Commands*}
window similar to Emacs's own \texttt{*scratch*}: type an
S-expression, press \texttt{C-j}, and the result is sent to the
running prover and inserted below.  \texttt{TAB} completes proof
command names; \texttt{C-c C-d} describes the command at point.

\paragraph{Programmatic evaluation.}
\texttt{(vnb-eval-string STR \&optional TIMEOUT)} sends \texttt{STR}
to the running prover, waits up to \texttt{TIMEOUT} seconds
(default~30) for the next REPL prompt, and returns the Scheme
return value as an Emacs string.  Use it from your own Elisp when
integrating the prover into a larger setup.

%% =========================================================
\section{A complete proof}
\label{sec:quickstart}

\begin{sloppypar}
Launch the workspace as in Section~\ref{sec:emacs} (\texttt{\textasciitilde/prover/VNB}).
The Home Workspace appears; the proof state below is what you would
see after stepping through the example in the Focus Workspace.  The
same commands typed at the bare REPL produce identical output.
\end{sloppypar}

We prove the classical syllogism: all men are mortal; Socrates is a
man; therefore Socrates is mortal.  There is no ``theory'' to extend
(Section~\ref{sec:structures-not-theories}): the premises are either
hypotheses of one statement, or the axioms of a structure.  Both
readings follow; run them in this order in one session, because the
second reserves the names it declares.

\paragraph{First reading: one statement, free variables.}
The classes \texttt{man} and \texttt{mortal} and the individual
\texttt{socrates} are free variables of the theorem; the premises are
its antecedent.

\begin{VNBsnip}
(sp (make-wff-from-string
  "(forall([x in man], x in mortal) and socrates in man)
     implies socrates in mortal"))
(di)
(ai 1)
(inst "forall([x in man], x in mortal)" "socrates")
(bc   "socrates in man implies socrates in mortal")
(ass)
(qed  'socrates-is-mortal)
\end{VNBsnip}

Final \texttt{*VNB State*}: \texttt{Proof complete.}  After \texttt{(di)}
the antecedent is assumption~1; \texttt{(ai 1)} splits the conjunction
into two assumptions; \texttt{inst} lands the implication for
\texttt{socrates}; \texttt{bc} turns the goal into its antecedent, which
\texttt{ass} finds in the context.

\paragraph{Second reading: a structure.}
Declare the species whose models are the worlds in which the syllogism
holds: a carrier of people, a carrier of mortals, a distinguished person,
and one law.  The theorem then quantifies over every model.

\begin{VNBsnip}
(declare-structure HUMANITY
  (instance-var s)
  (carriers PEOPLE MORTALS)
  (constant SOCRATES PEOPLE)
  (law "forall([x in PEOPLE(s)], x in MORTALS(s))"))

(sp (make-wff-from-string
  "forall([s in HUMANITY], SOCRATES(s) in MORTALS(s))"))
(di)
(mac-h 'humanity-class 1)
(mac-h 'is-humanity 1)
(dk-split-all!)
(inst "forall([x in PEOPLE(s)], x in MORTALS(s))" "SOCRATES(s)")
(bc   "SOCRATES(s) in PEOPLE(s) implies SOCRATES(s) in MORTALS(s)")
(ass)
(qed  'socrates-is-mortal-in-humanity)
\end{VNBsnip}

\texttt{(di)} lands \texttt{s in humanity}; the class axiom
\texttt{humanity-class} rewrites it to \texttt{is-humanity(s)}, and the
definition \texttt{is-humanity} unfolds that to the structure's typing
conditions and its law, one conjunction, which \texttt{dk-split-all!}
splits into separate assumptions; the rest is the first reading with
\texttt{socrates(s)}, \texttt{people(s)}, \texttt{mortals(s)} in place of
the free variables.

\begin{remark}
The same proofs in the styled workspace: \textbf{Start Proof} with the
statement, then the keys of the Focus Workspace (\texttt{d} for
\texttt{di}, \texttt{a} for \texttt{ass}, \ldots); a command with an
argument asks for it in the minibuffer.  The structure declaration is typed once in the
\texttt{*VNB Commands*} scratch sheet, or put in a file loaded at start.
The error panel (Section~\ref{sec:emacs}) catches typos as you go.
\end{remark}

\paragraph{What each command does.}
\begin{sloppypar}
\begin{itemize}
\item \texttt{(sp w)} starts a new proof with goal \texttt{w}.
\item \texttt{(di)} direct inference: introduces the goal's leading
  universals and the antecedent of an implication as assumptions.
\item \texttt{(ai f)} antecedent inference on assumption \texttt{f}
  (a formula or its 1-based index): a conjunction is split, an
  implication whose antecedent is in context is detached.
\item \texttt{(mac-h 'name k)} rewrites assumption \texttt{k} by the
  named definition or theorem; \texttt{(dk-split-all!)} splits every
  conjunction in the context.
\item \texttt{(inst f t)} instantiates the universally quantified
  assumption \texttt{f} at term \texttt{t}.
\item \texttt{(bc f)} backchains: assumption \texttt{f} is $P
  \Rightarrow Q$, current goal is $Q$, new goal is $P$.
\item \texttt{(ass)} closes a goal that appears in the assumption
  list.
\item \texttt{(qed 'name)} installs the completed proof as a named
  theorem.
\end{itemize}
\end{sloppypar}

\begin{remark}
The ``open goals'' count in \texttt{*VNB State*} climbs as the proof
progresses: each inference adds a node to the deduction graph, and
the parent stays open until its child grounds.  Completion is the
root grounding, not the count returning to zero in entry order.
Keep going until \texttt{Proof complete.} appears.
\end{remark}

\paragraph{The same proof, scripted.}
For batch use, the first reading in the long forms:
\begin{VNBsnip}
(let* ((wic (make-wff-from-string
              "(forall([x in man], x in mortal) and socrates in man)
                 implies socrates in mortal"))
       (ps  (start-proof wic))
       (ps  (cmd-direct-inference ps))
       (ps  (cmd-antecedent-inference ps
              (parse-string
                "forall([x in man], x in mortal) and socrates in man")))
       (ps  (cmd-instantiate ps
              (parse-string "forall([x in man], x in mortal)")
              'socrates))
       (ps  (cmd-backchain ps
              (parse-string
                "socrates in man implies socrates in mortal")))
       (ps  (cmd-assumption ps)))
  (cmd-qed ps 'socrates-is-mortal))
\end{VNBsnip}
All examples in this manual use the interactive short forms; the long
form is available for scripting.

\paragraph{What \texttt{make-wff} checks.}
\texttt{make-wff} (and \texttt{make-wff-from-string}, which calls it)
performs a position-typed walk of its argument and signals an error
rather than wrap a malformed expression.  The checks are:
\begin{sloppypar}
\begin{itemize}
\item \textbf{Numbers in wff positions} are rejected:
  \texttt{(make-wff-from-string "1 and 35")} $\Rightarrow$ ``number in wff
  position 1''.
\item \textbf{Term-forming heads in wff positions} are rejected:
  \texttt{(make-wff-from-string "choice(A)")} as a top-level wff fails because
  \texttt{choice} produces a term, not a wff.
\item \textbf{Wff-forming heads in term positions} are rejected:
  \texttt{(make-wff-from-string "(A and B) in C")} fails because \texttt{A and B}
  is a wff and the first argument of \texttt{in} must be a term.
\item \textbf{Bare symbols in wff positions} are rejected.  Only
  \texttt{truth} and \texttt{falsity} are valid bare-symbol wffs;
  anything else (e.g.\ \texttt{A implies B}) is an error.
\item \textbf{Free-symbol role conflicts}: a symbol that appears free
  in both wff and term positions in the same formula is rejected.
\item \textbf{Wff-position binders}: a binder whose body uses the
  bound variable in wff position is rejected within its scope.  This
  catches \texttt{forall([p], p implies p)}: \texttt{p} is bound as a
  class variable but the body uses it as a proposition.
\item \textbf{Bound vs.\ free same-name use} is \emph{accepted with a
  warning}.  In \texttt{forall([a], a in nn) implies a in zz} the
  inner \texttt{a} is bound and the outer \texttt{a} is free.  The
  validator emits a \texttt{;; warning} but does not reject.
\item \textbf{Arities and bound-variable shapes} are checked: every
  connective and quantifier gets the right number of arguments;
  \texttt{forall}, \texttt{forsome}, \texttt{iota}, \texttt{comp},
  \texttt{sep}, \texttt{big-union} and \texttt{vnb-lambda} bound variables must be
  symbols (\texttt{validate-wff!}, \texttt{wff.scm}); the parser requires each
  binding of a \texttt{lambdoid} to have the form \texttt{x in A}, and
  rejects \texttt{lambda}, which is not a binder.
\end{itemize}
\end{sloppypar}
%
\begin{remark}\label{remark:27135-34761-138906}
%
The set-theoretic \texttt{power(A)} (power set) and arithmetic
\texttt{x \^{} n} (exponentiation) share the same underlying symbol because
MIT~Scheme reads identifiers case-insensitively; \texttt{make-wff}
distinguishes them by arity (1 vs.\ 2 arguments).  All downstream
operations (\texttt{free-vars}, \texttt{subst-free},
\texttt{alpha-equiv?}, macete pattern matching) honour the same
arity-based dispatch.

%
\end{remark}

%% =========================================================
\section{A macete example}
\label{sec:quickstart-macete}

Macetes are named rewrite rules derived from installed theorems.
To apply a macete outside a proof, install a theorem and call
\texttt{apply-macete}:

\begin{VNBcode}
(install-theorem! 'eq-symm
  (make-wff-from-string "forall([x, y], x = y iff y = x)"))

(apply-macete 'eq-symm (make-wff-from-string "3 + 5 = 8"))
; => wff for "8 = 3 + 5"

(apply-macete 'eq-symm (make-wff-from-string "n = m implies truth"))
; => wff for "m = n implies truth"

(apply-macete 'eq-symm (make-wff-from-string "3 in nn"))
; => #f          -- pattern x = y does not match; no rewrite fires
\end{VNBcode}

\texttt{apply-macete} accepts and returns a \texttt{wff} object; it
returns \texttt{\#f} when the pattern matches nowhere.
To apply a macete as a proof step, use
\texttt{(mac \textit{name})} inside a proof (see
Section~\ref{sec:macetes}).

\begin{remark}
To try the four lines above in the styled workspace, open the
\texttt{*VNB Commands*} scratch sheet (\texttt{M-x vnb-command-buffer})
and press \texttt{C-j} after each S-expression; the prover's reply
appears on the line below.  The Home, Focus, and Proof Overview
workspaces are for \emph{proofs}; ad-hoc expression evaluation lives
on the scratch sheet or in the bare REPL.
\end{remark}

%% =========================================================


%% =========================================================
